\subsection{Death-receptor models and the commitment question}\label{sec:death-models}

We import the original 28-node Boolean GINsim model distributed with the
Calzone study, rather than a later multiscale adaptation.\cite{Calzone2010}
DR-FB+ uses its published rules. DR-FB- changes only the CASP8 rule from
$(\mathrm{DISC\_TNF}\lor\mathrm{DISC\_FAS}\lor\mathrm{CASP3})\land
\neg\mathrm{cFLIP}$ to
$(\mathrm{DISC\_TNF}\lor\mathrm{DISC\_FAS})\land\neg\mathrm{cFLIP}$.
CASP3 remains a dynamic component; this is not its knockout.
The source article's Figures S2 and S3 explicitly examine removal of this
interaction under sustained and transient stimulation. Its previously reported
biological interpretation is therefore a rationale for the comparison, not
independent new experimental validation of our results.

The common initial state has TNF, FADD, ATP and cIAP equal to one, and all other
nodes equal to zero. FASL is zero. Each asynchronous update changes one enabled
component; FASL and FADD remain fixed. The baseline protocol also fixes TNF at
one. The intervention protocol additionally enables \texttt{withdraw\_TNF}
at every state with TNF equal to one; this action changes only TNF to zero,
irreversibly, with identical availability in both variants.

The single predeclared interface observes activation and deactivation of
NFkB, CASP3 and MPT, together with the withdrawal action. All other updates
are silent. These are the source model's survival, apoptotic and necrotic
readouts; ATP depletion is additionally required when classifying a terminal
state as necrotic. An apoptotic terminal signature has CASP3 active and
NFkB/MPT inactive; survival has only NFkB active; necrosis has only MPT active
and ATP depleted. The naive signature has all three markers inactive and ATP
present. Mixed or varying terminal signatures are reported as other, not
forced into a fate class. No such additional terminal class occurred here.

We compute reachable terminal strongly connected components and separately
check whether CASP3 can ever become inactive. A singleton reachable terminal
fate is not, in general, an inevitability claim under arbitrary scheduling.
We use \emph{model-level commitment under the declared withdrawal semantics}
for the selected state's unique apoptotic terminal fate together with
CASP3 persistence after immediate withdrawal. The selected continuation graphs
are acyclic, so their terminating post-withdrawal behavior does not require an
extra fairness assumption. This does not establish cellular commitment or
recovery from an executed biological death program.

The initial object-based enumeration exceeded 200,000 states. Before relation
results were available, we documented an unchanged-interface implementation
amendment: omit updates of the hidden sink readouts NonACD, Apoptosis and
Survival, which regulate no retained node, and enumerate retained states using
compact Boolean/CSR storage. States agreeing on all retained coordinates form
a weak bisimulation between the full model and this quotient: a sink update
stutters, while every retained update and withdrawal lifts in either direction.
All rules and initial values for the remaining components are unchanged.
Weak traces, weak relations and retained terminal signatures are preserved;
strong bisimilarity of a quotient is not automatically a strong claim about
the unreduced model. Resource amendments, source hashes and the original
configuration are retained in the reproducibility record.

